\section{Tepeyac After Conquest: Place, Wound, and Mission}

The received date lies ten years after the fall of Mexico-Tenochtitlan in 1521. Spanish military victory did not create a blank religious landscape. Nahua peoples possessed cities, kinship orders, languages, political rivalries, ritual systems, schools, poetry, tribute networks, and memories of catastrophic war. European invasion was followed by epidemic mortality and famine, tribute and coerced labor, displacement, new institutions, and a contested missionary encounter. Indigenous rulers and communities also acted: they allied and resisted, litigated, recorded catastrophe in pictorial and alphabetic histories, negotiated labor burdens, and appropriated Christian institutions within conditions they did not choose.

No Catholic account of Guadalupe should turn conquest into a sacrament. The Gospel judges avarice, coercion, racial contempt, sexual violence, enslavement, and the destruction of persons. Neither can mission be reduced to conquest. Indigenous Christians truly appropriated the faith, recorded their communities, litigated, reshaped Christian institutions, entered sacramental life, and formed local churches that cannot be described as merely European religion imposed upon passive victims. Guadalupe's power arises within this morally mixed history.

\subsection{A Christian messenger with an Indigenous name}

The \emph{Nican mopohua} presents Juan Diego as a baptized, practicing Christian traveling to Tlatelolco to pursue worship and instruction. He is not a noble or colonial official. He speaks deferentially, recognizes episcopal authority, cares for his sick uncle, perseveres through rejection, and becomes a messenger whose credibility must be heard by a bishop.

His double name matters. “Juan Diego” marks baptismal incorporation and a Christian patron; “Cuauhtlatoatzin,” used in modern papal teaching and canonization, preserves Indigenous personhood. Baptism does not require cultural nonexistence. John Paul II made this point explicit in 2002: Juan Diego received the Christian message without renouncing his Indigenous identity and thereby witnessed to a humanity in which peoples become God's children.

Modern readers must also resist using diminutive speech to infantilize him. The narrative's tender forms of address belong to Nahuatl rhetoric and a maternal relation; they do not make an adult Indigenous man a permanent child before European authority. The story reverses colonial expectations: the one initially least credited carries the word and sign by which the pastor learns.

\subsection{Tepeyac as a layered sacred landscape}

Tepeyac stood north of the island-city and along important routes. In the polemical appendix on “superstitions” to Book XI, Bernardino de Sahagún calls the place Tepeacac or Tepeaquilla, describes a former temple and pilgrimage associated with \emph{Tonantzin}, “our revered mother,” and complains that pilgrims also called the later church of Our Lady of Guadalupe by that name. His account is a colonial missionary intervention, not neutral ethnography; yet it is primary evidence for layered sacred geography, long-distance pilgrimage, and contested vocabulary. The surviving evidence does not justify a simple equation “Guadalupe equals an unchanged goddess,” nor the opposite fantasy that the place possessed no prior sacred memory.

Christian inculturation neither baptizes every inherited cult nor destroys every cultural form. It discerns. Words, flowers, song, kinship address, pilgrimage, sacred landscape, and visual symbols can be assumed, purified, and redirected toward the triune God. Claims incompatible with Christ are rejected; genuine human goods are healed and fulfilled. The Marian title is Catholic because the woman is identified as the ever-Virgin Mother of the true God and sends her messenger into ecclesial communion. The narrative is Indigenous in language and imagination without becoming a covert second religion.

\begin{threestable}{Historical reality}{Theological responsibility}{Guadalupan implication}
Military conquest and colonial power & Name violence, coercion, exploitation, and unequal archives without treating every missionary or convert as one thing & Mary cannot be made patroness of domination; the poor messenger's dignity judges both Church and state.\\
Indigenous agency and conversion & Recognize baptized Nahuas as subjects of faith, thought, language, and mission & Juan Diego is an evangelizing agent, not a decorative recipient of European religion.\\
Earlier sacred geography & Distinguish continuity of place and vocabulary from identity of divine referent & Tepeyac can become a Christian place through discernment and conversion, not by semantic sleight of hand.\\
Ecclesial office & A bishop must test a claim and govern public worship & The messenger challenges the pastor to listen; the pastor's reception gives the sign a house in the Church.\\
\end{threestable}

\subsection{The “sacred little house”}

The narrative's central request is not a throne for ethnic or clerical power but a \emph{teocalli} or sacred house where Mary will manifest the one true God and give maternal compassion, help, and protection to the peoples who seek her. Modern Popes repeatedly return to this request. The house is at once architectural, ecclesial, and social.

Architecturally, the shrine makes memory inhabitable. Ecclesially, it becomes a place where Word, sacrament, episcopal oversight, popular piety, and pilgrimage meet. Socially, it asks whether those marginalized by conquest can enter not as tolerated strangers but as speakers and children. The Basilica can contradict its own origin if commerce eclipses prayer, Indigenous persons become imagery without voice, or the poor cannot find welcome.

The request also protects Marian theology. Mary asks for a house in order to show the true God, not to replace him. Her motherhood is missionary because it gathers a people toward Christ. John Paul II therefore describes genuine Marian devotion as a sure path to encounter with Christ; Leo XIV concludes the Guadalupan message with Mary's own Gospel direction: do whatever her Son tells you.

\subsection{Evangelization without numerical mythology}

Guadalupe is often credited with a sudden conversion of eight or nine million Indigenous persons within a few years. The numbers circulate without an identified baptismal census capable of sustaining that precision, and the foundational 1648--1649 texts do not supply the familiar apologetic total. Massive sixteenth-century evangelization is historical; its pace, motives, sacramental preparation, and relation to epidemic and social collapse require careful reconstruction.

Theologically, numerical exaggeration weakens rather than strengthens mission. Baptism is not a conquest statistic. Conversion can contain conviction, communal pressure, political adaptation, fear, grace, mixed motives, incomplete catechesis, and later deepening. Guadalupe's evangelizing fruit is better measured by the durable formation of Christian peoples, liturgy, charity, vocation, solidarity, and the capacity of the Gospel to take flesh in living languages.

\sourceboundary{The colonial setting and the received narrative's baptized Nahua characters}{The event's historical world: conquest, epidemic, mission, unequal authority, Indigenous agency, and a Church being formed across cultures}{A proof that colonial rule was holy, a census of conversions caused by the image, or a single formula that resolves the encounter between Christianity and every Indigenous tradition}
